% Ampliación acumulativa del TRIT; se inserta antes de la figura de regímenes.
% No sustituye ni elimina el desarrollo precedente.
\subsubsection{Division équilibrée et composition des relèvements}
\label{trit-ampl-div-balanceada}

La réduction ternaire permet d'exprimer une évaluation entière d'APP en une
coordonnée visible et une coordonnée de prolongation. Sa fonction ne consiste
pas seulement à abréger l'écriture d'un entier : elle fournit une règle de
composition qui détermine comment le quotient est modifié lorsqu'on additionne ou
multiplie des évaluations. Le choix équilibré explicite en outre la
compatibilité de cette règle avec l'inversion du signe.

\begin{definition}[Coordonnées ternaires équilibrées]
\label{trit-ampl-def-coordenadas}
Pour \(n\in\mathbb Z\), soient
\[
 q_3(n)=\left\lfloor\frac{n+1}{3}\right\rfloor,\qquad
 r_3(n)=n-3q_3(n).
\]
L'application
\[
 \mathcal B:\mathbb Z\longrightarrow\TRIT\times\mathbb Z,\qquad
 n\longmapsto(r_3(n),q_3(n))
\]
est appelée décomposition ternaire équilibrée. Son inverse est
\(\mathcal L(r,c)=r+3c\).
\end{definition}

En effet, \(r_3(n)\in\{-1,0,+1\}\). Si deux paires représentent le même entier,
\(r+3c=s+3d\), alors \(r-s\) est un multiple de \(3\) et appartient à l'intervalle
\([-2,2]\) ; nécessairement \(r=s\) et \(c=d\). La représentation est donc
unique. La division peut être répétée sur \(q_3(n)\), de sorte que chaque niveau
conserve la donnée nécessaire pour reconstruire le précédent. Le quotient n'est pas une
perturbation du résidu : il constitue sa coordonnée complémentaire.

La signature obtenue à partir d'une évaluation \(F\), avec \(F=S\) ou \(F=P\) dans APP, est \(r_3\circ F\). Lorsque la réduction agit sur le curseur de phase, la même section équilibrée produit les classes \(\{1,4,7\}\), \(\{2,5,8\}\) et \(\{3,6,9\}\), de valeurs \(+1,-1,0\). Le sélecteur du TPK défini ci-après utilise ces classes pour déterminer le mode opératoire. La lecture d'une évaluation et la sélection d'un mode ont ainsi la même réduction ternaire, mais des domaines et des fonctions différents.

\Needspace{16\baselineskip}
\begin{proposition}[Arithmétique des couples équilibrés]
\label{trit-ampl-prop-aritmetica}
Soient \(n=r+3c\) et \(m=s+3d\), avec \(r,s\in\TRIT\). Définissons
\[
 \chi(r,s)=\frac{r+s-r_3(r+s)}3.
\]
Les opérations entières sont transportées vers les couples au moyen de
\begin{align}
 (r,c)\boxplus(s,d)
 &=\bigl(r_3(r+s),c+d+\chi(r,s)\bigr),
 \label{trit-ampl-eq-suma}\\
 (r,c)\boxtimes(s,d)
 &=\bigl(rs,rd+sc+3cd\bigr).
 \label{trit-ampl-eq-producto}
\end{align}
L'application \(\mathcal L\) conserve les deux opérations.
\end{proposition}
\begin{proof}
L'égalité \(r+s=r_3(r+s)+3\chi(r,s)\) donne
\[
 n+m=r_3(r+s)+3\bigl(c+d+\chi(r,s)\bigr).
\]
D'autre part,
\[
 nm=rs+3(rd+sc+3cd).
\]
Comme le produit de deux représentants équilibrés appartient de nouveau à
\(\{-1,0,+1\}\), il ne nécessite pas de réduction résiduelle supplémentaire. L'unicité
de la décomposition démontre les deux formules.
\end{proof}

Par exemple, \(2=(-1)+3\) et \(4=1+3\). Leur somme et leur produit se reconstruisent
sous la forme
\[
 (-1,1)\boxplus(1,1)=(0,2),\qquad
 (-1,1)\boxtimes(1,1)=(-1,3).
\]
Les évaluations résultantes sont \(6\) et \(8\). La somme des quotients suffit
dans cet exemple additif ; dans le produit interviennent les termes croisés
\(rd\), \(sc\) et \(3cd\). L'écart entre les deux opérations réside
également dans la coordonnée de prolongement, et pas seulement dans le résidu.

Le terme \(\chi\) est un cocycle de la section équilibrée. Si
\(r\oplus s=r_3(r+s)\), il satisfait
\[
 \chi(r,s)+\chi(r\oplus s,t)
 =\chi(s,t)+\chi(r,s\oplus t).
\]
Les deux membres représentent le quotient de la même somme de trois entiers.
Cette identité garantit que le regroupement des opérations ne change pas l'
évaluation reconstruite. En même temps, la projection sur la première
composante supprime cette comptabilité : \(1+1=-1+3\), tandis que la lecture
résiduelle conserve uniquement \(-1\). L'information éliminée est
déterminée, dans ce cas, par \(\chi(1,1)=1\).

La relation avec le module \(9\) conserve une seconde échelle discrète. L'application
\[
 \beta:\TRIT\times\mathbb F_3\longrightarrow\mathbb Z/9\mathbb Z,
 \qquad (r,\bar c)\longmapsto r+3c\pmod9
\]
est bien définie et bijective. Changer le représentant de \(\bar c\) ajoute un multiple de \(9\) ; réciproquement, une égalité d'images fixe d'abord \(r\) par réduction modulo \(3\), puis \(\bar c\). L'addition transportée par cette bijection contient le terme \(\chi(r,s)\) de \eqref{trit-ampl-eq-suma}, maintenant réduit modulo \(3\). Ainsi, les deux niveaux ternaires ne se composent pas par une addition indépendante des coordonnées : la normalisation du premier niveau modifie le second.

En particulier, les classes \(3\), \(6\) et \(0\) de
\(\mathbb Z/9\mathbb Z\) ont pour première composante \(0\), mais pour seconde
composante \(1\), \(2\) et \(0\), respectivement. La distinction déjà observée
entre la réduction directe et l'extraction de la coordonnée de l'idéal apparaît
comme une propriété de \(\beta\). Conserver cette coordonnée empêche que le
secteur résiduel nul se transforme artificiellement en une seule position.
Le module \(1000\), utilisé ensuite pour organiser les blocs décimaux,
remplit une autre fonction de représentation positionnelle ; son quotient est conservé
au moyen de la division correspondante. Un changement de base conserve l'entier
lorsqu'il retient le résidu et le quotient, tandis que la coïncidence des résidus
dans des bases différentes n'identifie pas à elle seule leurs états de transport.

\Needspace{11\baselineskip}
\subsubsection{Inversion, classe résiduelle et coordonnées non visibles}
\label{trit-ampl-inversion}

L'inversion additive se relève sans ambiguïté :
\[
 \mathcal B(-n)=(-r_3(n),-q_3(n)).
\]
En particulier, \(3\), \(0\) et \(-3\) ont pour coordonnées
\((0,1)\), \((0,0)\) et \((0,-1)\). Ils partagent une signature nulle, mais représentent des
évaluations distinctes. Dans une transition restreinte à un quotient d'
amplitude maximale \(1\), ce sont les trois possibilités de la fibre visible au-dessus de
zéro ; dans un relèvement entier général, cette fibre contient tous les couples
\((0,c)\), avec \(c\in\mathbb Z\).

La coordonnée de quotient récupère l'entier évalué, mais ne récupère pas à elle seule la trajectoire qui l'a produit. Deux chemins d'APP peuvent partager évaluation, résidu et quotient, et différer par l'orientation, la succession des modes ou l'enroulement. La décomposition \(\mathcal B\) est donc intégrée à l'état de transport, au lieu de le remplacer. La mémoire arithmétique d'une division et la mémoire d'une trajectoire sont des coordonnées liées, et non des synonymes.

Pour une suite de signatures de longueur \(n\), l'inversion orientée est
\[
 C_{\rm rev}(\tau_1,\ldots,\tau_n)
   =(-\tau_n,\ldots,-\tau_1).
\]
L'appliquer deux fois restitue la suite. L'opération inverse l'ordre des
positions et le signe de chaque signature ; elle conserve sa longueur et les
positions nulles dans l'ordre inversé. Dans le paramétrage adopté,
elle échange le retour avec mémoire, associé à \(+1\), et l'ouverture
bidirectionnelle, associée à \(-1\), tout en maintenant le seuil \(0\). Le transport
des autres données d'une trajectoire est spécifié au moyen de
l'involution correspondante du TPK.

Cette inversion de signatures doit être distinguée de la conjugaison dans une
algèbre quadratique. La première peut changer le type \(\tau\) ; la seconde,
construite ci-dessous, agit à l'intérieur du type déjà fixé. La
distinction empêche d'attribuer à une inversion de signe une équivalence
algébrique entre des régimes différents.

\subsubsection{Produit quadratique, conjugaison et norme algébrique}
\label{trit-ampl-norma}

Un élément de \(\mathbb A_\tau\) possède une expression unique \(aI+bJ_\tau\). La relation \(J_\tau^2=-\tau I\) détermine son produit :
\begin{equation}
 (aI+bJ_\tau)(cI+dJ_\tau)
   =(ac-\tau bd)I+(ad+bc)J_\tau.
 \label{trit-ampl-producto-cuadratico}
\end{equation}
Il s'agit d'une réalisation réelle du type déjà sélectionné. Les coordonnées \(a,b\) ne sont pas des résidus d'APP, et le produit \eqref{trit-ampl-producto-cuadratico} ne remplace pas \eqref{trit-ampl-eq-producto} ; il exprime une autre opération, dans le codomaine quadratique déclaré.

\begin{definition}[Conjugaison et norme algébrique]
\label{trit-ampl-def-norma}
La conjugaison du type \(\tau\) et sa norme algébrique sont
\[
 \overline{aI+bJ_\tau}=aI-bJ_\tau,\qquad
 N_\tau(aI+bJ_\tau)=a^2+\tau b^2.
\]
\end{definition}

\begin{proposition}[Multiplicativité et inversibilité]
\label{trit-ampl-prop-norma}
Pour \(z,w\in\mathbb A_\tau\),
\[
 z\bar z=N_\tau(z)I,\qquad
 N_\tau(zw)=N_\tau(z)N_\tau(w).
\]
De plus, \(z\) est inversible si et seulement si \(N_\tau(z)\ne0\), et dans ce cas
\[
 z^{-1}=\frac{\bar z}{N_\tau(z)}.
\]
\end{proposition}
\begin{proof}
Le produit de \(aI+bJ_\tau\) et de \(aI-bJ_\tau\) est
\((a^2+\tau b^2)I\). La conjugaison est multiplicative et l'algèbre est
commutative, de sorte que
\((zw)\overline{zw}=(z\bar z)(w\bar w)\). Si la norme ne s'annule pas, la
formule indiquée fournit l'inverse. Réciproquement, si \(zw=I\), la
multiplicativité implique \(N_\tau(z)N_\tau(w)=1\).
\end{proof}

La forme quadratique \(N_\tau\) n'est définie positive que dans le type elliptique. Pour \(\tau=0\), elle vaut \(a^2\), et l'élément \(J_0\) est non nul bien que son carré et sa norme soient nuls. Pour \(\tau=-1\), elle vaut \(a^2-b^2\) ; les éléments non nuls \(I+J_{-1}\) et \(I-J_{-1}\) ont une norme nulle et leur produit s'annule. La norme algébrique est ici une forme quadratique multiplicative : elle ne présuppose pas une norme positive d'espace vectoriel.

Les trois types peuvent être réalisés fidèlement au moyen de
\[
 aI+bJ_\tau\longmapsto
 \begin{pmatrix}a&-\tau b\\b&a\end{pmatrix}.
\]
Le déterminant de cette matrice est \(N_\tau(aI+bJ_\tau)\). Sont ainsi
représentées, avec une même expression matricielle, la structure
complexe, la structure des nombres duaux et la structure réelle scindée.
Chacune conserve sa loi et ses diviseurs de zéro. La forme de signature
\((1,1)\) du dernier cas a une interprétation lorentzienne en dimension
\(1+1\) ; l'identification d'une métrique physique exige en outre un espace,
un champ et une règle de transport spécifiques.

\subsubsection{Conservation multiplicative et composition des évolutions}
\label{trit-ampl-evoluciones}

L'exponentielle déjà obtenue appartient à l'ensemble des éléments de norme un :
\[
 N_\tau\bigl(\exp(tJ_\tau)\bigr)=1.
\]
Cette conservation se maintient lors de la composition des paramètres. En effet, tous les
facteurs sont des fonctions du même générateur, et
\[
 \exp(tJ_\tau)\exp(sJ_\tau)=\exp((t+s)J_\tau).
\]
En comparant les deux composantes, on déduit les lois d'addition
\begin{align}
 C_\tau(t+s)&=C_\tau(t)C_\tau(s)-\tau S_\tau(t)S_\tau(s),\\
 S_\tau(t+s)&=S_\tau(t)C_\tau(s)+C_\tau(t)S_\tau(s).
\end{align}
La conjugaison remplace \(t\) par \(-t\) et coïncide avec l'inversion dans
ce sous-groupe. Elle ne change pas \(\tau\) : elle inverse une évolution du régime
choisi.

Dans le type elliptique, la forme conservée est \(a^2+b^2\), et la trajectoire exponentielle parcourt le cercle de norme un. Dans le type hyperbolique, les coordonnées propres sont \(u=a+b\), \(v=a-b\), avec \(uv=N_{-1}(z)\). L'exponentielle donne
\[
 (u,v)=(e^t,e^{-t}),\qquad uv=1.
\]
La croissance d'une coordonnée s'accompagne de la contraction de l'autre.
La conservation exprime une relation entre les deux coordonnées ; elle n'affirme pas que
chacune reste constante.

Dans le type parabolique,
\[
 (I+tJ_0)(I+sJ_0)=I+(t+s)J_0,\qquad
 (I+tJ_0)^{-1}=I-tJ_0.
\]
L'évolution ne s'arrête pas du fait que \(J_0^2=0\). La nilpotence élimine les
termes d'ordre au moins deux, mais conserve le terme linéaire qui
transporte le paramètre. Le symbole visible \(0\), l'
élément nul de l'algèbre et le générateur non nul \(J_0\) sont distincts. Cette distinction
est nécessaire pour décrire le seuil sans l'identifier à l'absence d'
état ou de dynamique.

\begin{proposition}[Composition de coordonnées affines]
\label{trit-ampl-prop-pendientes}
Dans la coordinatisation où la composante scalaire ne s'annule pas, représentons une
direction par \(I+uJ_\tau\). Si \(1-\tau uv\ne0\), la direction de son
produit avec \(I+vJ_\tau\) a pour coordonnée
\[
 u\star_\tau v=\frac{u+v}{1-\tau uv}.
\]
\end{proposition}
\begin{proof}
On a
\[
 (I+uJ_\tau)(I+vJ_\tau)
   =(1-\tau uv)I+(u+v)J_\tau.
\]
La division par la composante scalaire donne l'expression annoncée.
\end{proof}

La coordinatisation affine conserve le rapport entre les composantes, et non le facteur scalaire
par lequel on a divisé. Lorsque \(1-\tau uv=0\), il faut examiner le produit avant de
changer de coordinatisation : il peut avoir une direction définie avec une composante scalaire
nulle, ou être l'élément nul en présence de diviseurs de zéro. La loi de
composition appartient toujours à l'algèbre complète ; la singularité d'une
coordonnée n'équivaut pas automatiquement à une singularité de l'objet.
Cette formule réapparaîtra dans la lecture régionale elliptique. Son emploi exige
de maintenir le dénominateur et, le cas échéant, le facteur de normalisation.

\subsubsection{Transport du type et réalisations ultérieures}
\label{trit-ampl-puentes}

Les algèbres \(\mathbb A_{+1}\), \(\mathbb A_0\) et
\(\mathbb A_{-1}\) ne sont pas isomorphes en tant qu'algèbres réelles unitaires.
La première est un corps ; la deuxième contient un nilpotent non nul ; la
troisième possède des idempotents non triviaux et est isomorphe à
\(\mathbb R\times\mathbb R\). Le changement de signature qui organise une
route ne peut donc pas s'interpréter comme un isomorphisme indifférencié entre les
trois réalisations. Le TPK conserve l'indice de régime, les domaines et
l'opérateur effectif qui agit à chaque transition.

\subsubsection{Norme du commutateur des projections arithmétiques}

On obtient le préfacteur électronique sans utiliser \(\alpha\), la torsion
ou une unité dimensionnelle. Le domaine est la troisième réalisation de APP,
\(\mathcal A_3=\{1,\ldots,9\}^3\), munie de la mesure uniforme.
L’espace de fonctions est \(\ell^2(\mathcal A_3;\mathbb R)\), muni du produit
scalaire
\[
 \langle f,g\rangle=\frac1{729}\sum_{x\in\mathcal A_3}f(x)g(x).
\]
Dans ce bloc, \([P,Q]=PQ-QP\) désigne le commutateur additif d’opérateurs.
Nous définissons
\[
 \Sigma(x)=\rho_9(x_1+x_2+x_3),\qquad
 \Pi(x)=\rho_9(x_1x_2x_3).
\]
Les espérances conditionnelles \(P_\Sigma\) et \(P_\Pi\) sont les
projections orthogonales sur les fonctions constantes sur chaque fibre
additive et multiplicative, respectivement. Ainsi, pour une fonction \(f\),
\[
 (P_\Sigma f)(x)
 =\frac{1}{|\Sigma^{-1}(\Sigma(x))|}
   \sum_{y:\,\Sigma(y)=\Sigma(x)}f(y),
\]
et de même pour \(P_\Pi\). Les deux projections agissent sur le même espace ; elles ne représentent pas deux supports arithmétiques indépendants.

\begin{lemma}[Deux projections et angles principaux]
Si \(s\in[0,1]\) est une valeur singulière de l'appariement entre des bases
orthonormales des images de deux projections \(P,Q\), son bloc non
trivial contribue à \(\|[P,Q]\|\) par
\(s\sqrt{1-s^2}\).
\end{lemma}
\begin{proof}
Dans le plan principal correspondant, on peut écrire
\[
 P=\begin{pmatrix}1&0\\0&0\end{pmatrix},\qquad
 Q=\begin{pmatrix}
 s^2&s\sqrt{1-s^2}\\
 s\sqrt{1-s^2}&1-s^2
 \end{pmatrix}.
\]
Leur commutateur est
\[
 [P,Q]=s\sqrt{1-s^2}
 \begin{pmatrix}0&1\\-1&0\end{pmatrix}.
\]
La matrice finale est orthogonale, donc sa norme est un. Les blocs communs ou orthogonaux ont un commutateur nul. La décomposition orthogonale en plans principaux donne l'affirmation et réduit la norme totale au maximum de ces contributions.
\end{proof}

\begin{proposition}[Norme exacte d'incompatibilité]
Sur \(\ell^2(\mathcal A_3)\),
\begin{equation}
 \|[P_\Sigma,P_\Pi]\|=\frac{\sqrt3}{4}.
 \label{ivloc:norma}
\end{equation}
\end{proposition}
\begin{proof}
Soit \(N_{ab}=|\Sigma^{-1}(a)\cap\Pi^{-1}(b)|\). L'énumération entière
des trois symboles fournit
\[
 N=9\begin{pmatrix}
 0&1&1&0&1&1&0&1&4\\
 1&0&1&1&0&1&1&0&4\\
 1&0&2&0&1&2&0&0&3\\
 0&1&1&0&1&1&0&1&4\\
 1&0&1&1&0&1&1&0&4\\
 0&0&2&1&0&2&0&1&3\\
 0&1&1&0&1&1&0&1&4\\
 1&0&1&1&0&1&1&0&4\\
 0&1&2&0&0&2&1&0&3
 \end{pmatrix}.
\]
Chaque ligne a pour somme \(81\) ; les sommes des colonnes sont
\[
 (m_1,\ldots,m_9)=(36,36,108,36,36,108,36,36,297).
\]
Les indicatrices normalisées des fibres ont pour matrice
d'appariement \(C_{ab}=N_{ab}/\sqrt{81m_b}\). Par conséquent,
\[
 M=CC^{\mathsf T},\qquad
 M_{ac}=\sum_{b=1}^9\frac{N_{ab}N_{cb}}{81m_b}
\]
est une matrice rationnelle dont le spectre se compose des carrés des
valeurs singulières requises.

La multiplication exacte montre que \((1,\ldots,1)\),
\((1,-1,0)^3\) et \((1,1,-2)^3\) sont des vecteurs propres de valeurs propres
\(1\), \(1/4\) et \(7/132\), respectivement. Le sous-espace des
vecteurs à support dans \(\{3,6,9\}\) dont la somme est nulle a pour dimension
deux et pour valeur propre \(1/18\). Les vecteurs de somme nulle à support dans
\(\{1,4,7\}\) ou dans \(\{2,5,8\}\) forment un noyau de dimension
quatre. Ces sous-espaces sont indépendants et la somme de leurs dimensions vaut
neuf. Le spectre complet est donc démontré
\[
 \operatorname{spec}(M)=
 \left\{1,\frac14,\frac7{132},\frac1{18},\frac1{18},0,0,0,0\right\}.
\]
Le maximum de \(\lambda(1-\lambda)\) sur cet ensemble est \(3/16\), atteint en \(\lambda=1/4\). Le lemme précédent donne \(\sqrt{3/16}=\sqrt3/4\).
\end{proof}

Le mode \((1,-1,0)^3\) est précisément la distribution des valeurs du
sélecteur tritique de phase dans la coordinatisation nonadique. La norme scalaire retient
une mesure d'incompatibilité, mais ne conserve pas à elle seule
l'orientation de ce mode et ne reconstruit pas les deux projecteurs. Cette mémoire
demeure dans l'état de route qui accompagne la lecture scalaire.

\begin{proposition}[Algèbre du bloc électronique local]
\label{ivloc:algebra}
Dans le plan principal qui atteint \eqref{ivloc:norma}, soient
\[
 P=\begin{pmatrix}1&0\\0&0\end{pmatrix},\qquad
 Q=\begin{pmatrix}1/4&\sqrt3/4\\\sqrt3/4&3/4\end{pmatrix},\qquad
 S=2P-I,\qquad J=\frac4{\sqrt3}[P,Q].
\]
Alors \(S^2=I\), \(J^2=-I\), \(SJ=-JS\), et l'algèbre réelle engendrée par \(S,J\) est \(M_2(\mathbb R)\), réalisation de \(\mathrm{Cl}_{1,1}\). En outre,
\[
 n_\pm=\frac{S\pm J}{2},\qquad
 n_\pm^2=0,\qquad n_+n_-+n_-n_+=I.
\]
\end{proposition}
\begin{proof}
On a \(S=\operatorname{diag}(1,-1)\) et
\(J=\bigl(\begin{smallmatrix}0&1\\-1&0\end{smallmatrix}\bigr)\).
La multiplication vérifie les trois relations. Les matrices \(I,S,J,SJ\) sont linéairement indépendantes et forment une base de \(M_2(\mathbb R)\), ce qui prouve l'affirmation sur l'algèbre. Enfin, \((S\pm J)^2=S^2+J^2\pm(SJ+JS)=0\), et la somme des produits croisés est \((S^2-J^2)/2=I\).
\end{proof}

Dans ce bloc coexistent une direction elliptique, une direction hyperbolique et deux
directions nulles paraboliques. Le résultat décrit une représentation
matricielle locale de l'arithmétique des deux feuillets. Il n'identifie pas toute APP
à \(M_2(\mathbb R)\), ni à un carré latin quantique ; de telles
affirmations exigeraient d'autres domaines et d'autres relations. L'exponentielle
n'appartient pas non plus exclusivement au régime parabolique : elle peut
s'appliquer à n'importe lequel des générateurs de ce bloc.


La proposition~\ref{ivloc:algebra} établit que la paire de projections des feuillets de APP
produit un commutateur normalisé dont le carré est \(-I\). Le même bloc
contient une involution de carré \(I\) et deux directions nilpotentes. La
réalisation conjointe est démontrée sur cette paire concrète de projections ;
la classification précédente explique pourquoi les trois relations quadratiques
peuvent apparaître dans une même algèbre matricielle sans être identifiées entre elles.

Dans le prolongement exceptionnel, la matrice de Paley réduite modulo \(3\) satisfait \(A_W^2=-I_6\). La relation dote l'espace ternaire d'une structure de dimension \(3\) sur \(\mathbb F_9\). Son lien avec une racine opératorielle réelle s'exprime au moyen de la présentation entière
\[
 \mathbb Z[u]/(u^2+1).
\]
Les changements de base vers \(\mathbb F_3\) et vers \(\mathbb R\) produisent des
représentations de cette même relation. Chacune conserve sa caractéristique,
sa dimension et son information d'incidence. Le développement des
opérateurs respectifs établit le lien, et non l'égalité de leurs
cardinaux.

La conservation de la norme un décrit donc les réalisations
quadratiques des évolutions. La conservation des quotients, de l'orientation et de la
trajectoire appartient à l'état arithmétique transporté. Les deux propriétés
interviennent dans la construction, avec des fonctions différentes : l'une détermine les
identités des représentations ; l'autre permet de composer, de reconstruire et de
prolonger les opérations qui les produisent.
